\chapter{Risque de crédit de portefeuille}\label{ch:rc:portfolio-credit}

Au début de 2005, une stratégie très prisée vendait une protection sur la tranche equity de l'indice
iTraxx Europe, soit les 3~\% premières pertes sur 125 noms, et la couvrait contre les mouvements
de spreads de marché en achetant une protection sur l'indice~; la Banque des règlements internationaux
évaluait son carry à 300 à 400 points de base par an. En mai, après les dégradations de General
Motors et de Ford (chapitre~14), la corrélation implicite dans les cours de la tranche equity a chuté
brutalement, le spread de la tranche equity s'est élargi, et les dénouements ont fait baisser les spreads
des tranches mezzanine, parfois dans le sens opposé à celui de l'indice. Les positions
couvertes par chacun des modèles de leurs détenteurs ont perdu sur les deux jambes au cours de
ces mêmes semaines. Ce chapitre construit ces modèles~: la dépendance des défauts par une copule, la
distribution des pertes d'un portefeuille, les prix et les risques des tranches, et pourquoi une corrélation
implicite d'une tranche est une convention de cotation dont les variations constituent un risque à part entière.
Le livre~2, chapitre~24, a introduit les tranches, le grand portefeuille et la corrélation de base~;
ici le portefeuille est fini, les tranches ont une structure par terme, et les couvertures sont mesurées.

\section{Corrélation de défaut et copules}

\begin{definition}[Corrélation de défaut]\label{def:rc:portfolio-credit:dc}
La \emph{corrélation de défaut}\index{corrélation de défaut} de deux noms sur un horizon est
la corrélation de leurs indicateurs de défaut $\mathbf 1_{\tau_i\le T}$ et
$\mathbf 1_{\tau_j\le T}$~:
\[
\rho^{D}_{ij} = \frac{\P(\tau_i\le T,\tau_j\le T)-p_ip_j}{\sqrt{p_i(1-p_i)p_j(1-p_j)}}.
\]
\end{definition}

Une copule (livre~4, chapitre~15) relie les lois marginales des instants de défaut,
chacune donnée par la courbe de survie d'un nom du chapitre~13, en une loi conjointe.
Le choix du marché est le plus simple.

\begin{definition}[Copule gaussienne]\label{def:rc:portfolio-credit:gauss}
Dans la \emph{copule gaussienne}\index{copule gaussienne} à un facteur, chaque nom possède une
variable latente $X_i = \sqrt\rho\,Z+\sqrt{1-\rho}\,\varepsilon_i$, où $Z,\varepsilon_i$ sont
des normales centrées réduites indépendantes, et fait défaut avant $t$ si $X_i < N^{-1}(p_i(t))$, où $p_i(t) = 1-Q_i(t)$.
Chaque loi marginale est respectée~; $\rho$, la corrélation des variables
latentes (la corrélation d'actifs), fixe la dépendance.
\end{definition}

La corrélation d'actifs n'est pas la \omterm{def:rc:portfolio-credit:dc}{corrélation de défaut}. Avec une \omterm{def:rc:reduced-form-credit:pd}{probabilité de défaut} à cinq ans
de 4.08~\% (un spread de 50 points de base, recouvrement de 40~\%), une corrélation d'actifs de 25~\% donne une
\omterm{def:rc:portfolio-credit:dc}{corrélation de défaut} de 7.0~\%~: deux défauts rares coïncident moins souvent que leurs
variables latentes évoluent ensemble.

\section{La copule gaussienne à un facteur et le grand portefeuille}

Conditionnellement au facteur commun $Z = z$, les noms sont indépendants, chacun faisant défaut avant $t$
avec la probabilité
\[
p_i(t\mid z) = N\!\left(\frac{N^{-1}(p_i(t))-\sqrt\rho\,z}{\sqrt{1-\rho}}\right).
\]
Pour un portefeuille fini, la loi du nombre de défauts sachant $z$ s'obtient en ajoutant les noms
un à un~: avec $P_k$ la probabilité de $k$ défauts parmi les premiers noms,
l'ajout du nom $i$ donne $P_k \leftarrow P_k(1-p_i)+P_{k-1}p_i$, la récurrence d'Andersen, Sidenius et
Basu (2003). En intégrant sur $z$ par quadrature, on obtient la loi des pertes inconditionnelle
(\cref{fig:rc:portfolio-credit:lossdist}).

\omcode{code/firm/portcredit/firm_portcredit.py}{40}{54}{La récurrence sur les noms conditionnellement au facteur, et la loi inconditionnelle du nombre de défauts par quadrature de Gauss--Hermite.}\label{lst:rc:portfolio-credit:recursion}

\begin{omfigure}
\begin{tikzpicture}
\begin{axis}[width=11.5cm,height=5.4cm,ymode=log,
  xlabel={nombre de défauts en cinq ans (125 noms)},ylabel={probabilité},
  tick label style={font=\small},label style={font=\small},
  xmin=0,xmax=40,ymin=1e-6,ymax=1,grid=major,grid style={gray!25},
  legend columns=3,legend style={font=\footnotesize,at={(0.5,-0.25)},anchor=north,draw=none,fill=none,/tikz/every even column/.append style={column sep=8pt}},
  table/col sep=comma]
\addplot[omMeth,very thick] table[x=k,y=r05]{figdata/rates-credit-risk/15-portfolio-credit/loss_dist.csv};
\addplot[omDef,very thick,dashed] table[x=k,y=r25]{figdata/rates-credit-risk/15-portfolio-credit/loss_dist.csv};
\addplot[omThm,very thick,dotted] table[x=k,y=r50]{figdata/rates-credit-risk/15-portfolio-credit/loss_dist.csv};
\legend{$\rho = 5~\%$,$\rho = 25~\%$,$\rho = 50~\%$}
\end{axis}
\end{tikzpicture}

\omcaption{Distribution du nombre de défauts en cinq ans dans un portefeuille de 125 noms
à 50 points de base, pour trois corrélations d'actifs (échelle logarithmique). La moyenne est la même, 5.1
défauts~; la corrélation déplace la probabilité à la fois vers zéro défaut et vers un grand nombre. Données~: le
tutoriel du chapitre.}\label{fig:rc:portfolio-credit:lossdist}
\end{omfigure}

\begin{definition}[Grand portefeuille homogène]\label{def:rc:portfolio-credit:lhp}
L'approximation du \emph{grand portefeuille homogène}\index{grand portefeuille homogène} (large homogeneous pool)
suppose une infinité de noms identiques~: sachant $z$, la fraction de perte vaut exactement
$(1-R)\,p(t\mid z)$, si bien que la perte du portefeuille est une fonction monotone du facteur et sa
loi est disponible sous forme fermée (la limite établie par Vasicek pour les portefeuilles de prêts).
\end{definition}

Le passage à la limite du grand portefeuille du livre~2 est rapide et proche pour de nombreux noms~: la perte espérée à cinq ans
de la tranche equity vaut 51.6~\% de son notionnel à la limite, contre 49.6~\% pour
125 noms à $\rho = 25~\%$. Le portefeuille fini compte pour la tranche equity et pour les portefeuilles
dont les noms diffèrent, ce à quoi sert la récurrence.

\section{Valorisation des tranches~: récurrence, corrélation composée et corrélation de base}

Une tranche $[a,d]$ perd $\min(\max(L_t-a,0),d-a)/(d-a)$ de son notionnel lorsque le portefeuille
a perdu $L_t$. Sa perte espérée à chaque date de prime, $\mathrm{EL}(t_k)$, donne les jambes comme pour
un credit default swap~: la prime est payée sur le notionnel restant espéré,
$\sum_k\delta_kP(0,t_k)\bigl(1-\tfrac12(\mathrm{EL}(t_{k-1})+\mathrm{EL}(t_k))\bigr)$, et la protection
verse $\sum_kP(0,t_k)\bigl(\mathrm{EL}(t_k)-\mathrm{EL}(t_{k-1})\bigr)$. La tranche equity se cote comme un
versement initial plus 500 points de base courus~; les autres se cotent en spread courant.

\begin{example}[Cotations synthétiques]\label{ex:rc:portfolio-credit:quotes}
Prenons 125 noms à 50 points de base, recouvrement de 40~\%, un taux plat de 4~\%, et un skew de
corrélation de base de 15~\%, 25~\%, 31~\%, 36~\% et 50~\% aux points de détachement de 3, 6, 9, 12 et 22~\% (à titre illustratif).
Les tranches à cinq ans se cotent à un versement initial de 37.21~\% plus 500 courus pour 0--3~\%, et
à 195, 83.6, 43.6, 17.9 et 3.6 points de base pour 3--6, 6--9, 9--12, 12--22 et 22--100~\%~;
l'indice est à 49.6.
\end{example}

\begin{definition}[Corrélation composée]\label{def:rc:portfolio-credit:compound}
La \emph{corrélation composée}\index{corrélation composée} d'une tranche est l'unique
corrélation pour laquelle la \omterm{def:rc:portfolio-credit:gauss}{copule gaussienne} à un facteur la reprice, avec ses pertes
d'attachement et de détachement calculées à cette corrélation.
\end{definition}

Les \omterm{def:rc:portfolio-credit:compound}{corrélations composées} des cotations synthétiques valent 15.0~\% pour la tranche equity, 12.9~\%,
18.4~\%, 25.2~\% et 50.0~\% pour les tranches au-dessus de 6~\%, et \emph{deux} valeurs pour la
mezzanine, 2.1~\% et 89.9~\%. Une mezzanine est longue en corrélation à faible corrélation (davantage de corrélation fait remonter les pertes
jusqu'à elle) et courte à forte corrélation (les pertes la dépassent d'un bond), donc
son spread est en bosse en $\rho$ (\cref{fig:rc:portfolio-credit:mezz}). Une cotation supérieure au sommet de la bosse,
soit 374 points de base, proche de $\rho = 29\%$, n'a aucune \omterm{def:rc:portfolio-credit:compound}{corrélation composée}. C'est pourquoi le marché
s'est tourné vers la corrélation de base du livre~2, définie sur les tranches de base
$[0,d]$ dont la valeur est monotone en corrélation~; une mezzanine $[a,d]$ se price alors à partir de deux
corrélations de base, $\mathrm{EL}_{a,d} = (d\,\mathrm{EL}_{0,d}(\rho_d)-a\,\mathrm{EL}_{0,a}(\rho_a))/(d-a)$.

\begin{omfigure}
\begin{tikzpicture}
\begin{axis}[width=11.5cm,height=5.4cm,
  xlabel={corrélation composée $\rho$},ylabel={spread par 3--6~\% (pb)},
  tick label style={font=\small},label style={font=\small},
  xmin=0,xmax=1,ymin=100,ymax=400,grid=major,grid style={gray!25},
  legend columns=2,legend style={font=\footnotesize,at={(0.5,-0.25)},anchor=north,draw=none,fill=none,/tikz/every even column/.append style={column sep=8pt}},
  table/col sep=comma]
\addplot[omDef,very thick] table[x=rho,y=par]{figdata/rates-credit-risk/15-portfolio-credit/mezz.csv};
\addplot[omThm,thick,dashed] table[x=rho,y=quote]{figdata/rates-credit-risk/15-portfolio-credit/mezz.csv};
\legend{spread modèle,cotation de marché 195 pb}
\end{axis}
\end{tikzpicture}

\omcaption{Le spread modèle de la tranche 3--6~\% en fonction de la \omterm{def:rc:portfolio-credit:compound}{corrélation composée}. La
courbe est en bosse, de sorte que la cotation de 195 points de base est atteinte en deux corrélations (2.1~\%
et 89.9~\%), et une cotation supérieure à la bosse n'en atteint aucune. Données~: tutoriel
du chapitre.}\label{fig:rc:portfolio-credit:mezz}
\end{omfigure}

\begin{definition}[Delta de tranche]\label{def:rc:portfolio-credit:delta}
Le \emph{delta de tranche}\index{delta de tranche} est le notionnel de protection sur l'indice dont la
valeur varie autant que celle de la tranche, par unité de notionnel de tranche, lorsque tous les spreads des noms
bougent ensemble d'une faible quantité, les corrélations de base étant maintenues.
\end{definition}

\begin{example}[Deltas sur l'ensemble de la structure]\label{ex:rc:portfolio-credit:deltas}
Les deltas des six tranches de l'\cref{ex:rc:portfolio-credit:quotes} valent 17.3, 7.32, 3.28,
1.77, 0.72 et 0.14 (\cref{fig:rc:portfolio-credit:deltas})~: 10~millions d'euros de tranche equity
bougent comme 173~millions d'euros d'indice. Couvrir l'equity avec la mezzanine demande
$17.3/7.32 = 2.36$ fois le notionnel.
\end{example}

\begin{omfigure}
\begin{tikzpicture}
\begin{axis}[width=11.5cm,height=5.2cm,ybar,bar width=16pt,
  symbolic x coords={0-3,3-6,6-9,9-12,12-22,22-100},xtick=data,
  xlabel={tranche (\%)},ylabel={delta (fois l'indice)},
  tick label style={font=\small},label style={font=\small},
  ymin=0,ymax=20,grid=major,grid style={gray!25},
  nodes near coords,every node near coord/.append style={font=\footnotesize},
  table/col sep=comma]
\addplot[fill=omDef!70,draw=omDef] table[x=tranche,y=delta]{figdata/rates-credit-risk/15-portfolio-credit/deltas.csv};
\end{axis}
\end{tikzpicture}

\omcaption{\omterm{def:rc:portfolio-credit:delta}{Deltas de tranche} face à un mouvement parallèle de l'indice, par unité de notionnel de tranche,
les corrélations de base étant maintenues. Le levier est maximal au bas de la structure~: la
tranche equity bouge comme 17~fois son notionnel d'indice, la super senior comme un septième. Données~: tutoriel
du chapitre.}\label{fig:rc:portfolio-credit:deltas}
\end{omfigure}

\begin{definition}[Tranche super senior]\label{def:rc:portfolio-credit:ss}
La \emph{tranche super senior}\index{tranche super senior} est le sommet d'une structure de capital,
au-dessus de la tranche qui obtiendrait la meilleure notation à elle seule (22--100~\% dans la
structure standard de l'indice européen)~: une perte attendue très faible sur un notionnel très grand,
la plus grande part de son risque se trouvant dans la valeur de marché de cette petite perte.
\end{definition}

\begin{remark}[La queue qui compte]\label{rem:rc:portfolio-credit:ss}
À une corrélation d'actifs de 50~\%, la probabilité que la perte à cinq ans du portefeuille dépasse 22~\% vaut
1.88~\%~: faible, sur le plus grand notionnel de la structure. En 2007--08, les \omterm{def:rc:portfolio-credit:ss}{tranches super senior}
de CDO adossés à des titres hypothécaires, conservées par des souscripteurs comme
Citigroup, Merrill Lynch et UBS ou assurées par AIG, ont causé des pertes de plusieurs milliards de dollars,
selon la Financial Crisis Inquiry Commission~: les détenteurs les avaient traitées comme
sans risque parce que leur perte attendue l'était. UBS seule a déprécié environ 21.7~milliards de dollars sur
les \omterm{def:rc:portfolio-credit:ss}{tranches super senior} qu'elle avait conservées, soit quelque 50~milliards de dollars avant la crise et sa plus grande source de
pertes~; le régulateur suisse a constaté que ces positions avaient été considérées comme neutres en VaR.
\end{remark}

\section{Au-delà de la gaussienne~: queues épaisses et dépendance de queue}

\begin{definition}[Copule de Student]\label{def:rc:portfolio-credit:t}
Dans la \emph{copule de Student}\index{copule de Student} à $\nu$ degrés de liberté, les
variables latentes sont $X_i = \sqrt W(\sqrt\rho\,Z+\sqrt{1-\rho}\,\varepsilon_i)$ avec $W = \nu/\chi^2_\nu$ commun
à tous les noms~; un nom fait défaut si $X_i < t_\nu^{-1}(p_i)$. L'échelle commune $W$
fait que les noms font défaut ensemble dans les mauvais tirages~: la copule présente une dépendance de queue, que la gaussienne
ne possède pas.
\end{definition}

\begin{example}[Queues]\label{ex:rc:portfolio-credit:tails}
À une corrélation d'actifs de 25~\%, la probabilité que le portefeuille perde plus de 12~\% en cinq
ans vaut 2.24~\% sous la \omterm{def:rc:portfolio-credit:gauss}{copule gaussienne} et 5.23~\% sous la \omterm{def:rc:portfolio-credit:t}{copule de Student} à quatre
degrés de liberté~; pour une perte supérieure à 24~\%, 0.13~\% contre 1.13~\%, soit 8.6~fois plus
(\cref{fig:rc:portfolio-credit:tails}). Les tranches seniors évaluées avec la \omterm{def:rc:portfolio-credit:gauss}{copule gaussienne}
à une seule corrélation paraissent bon marché pour la même raison que le skew existe~: le marché price plus de queue que
ne le donne la gaussienne.
\end{example}

\begin{omfigure}
\begin{tikzpicture}
\begin{axis}[width=11.5cm,height=5.4cm,ymode=log,
  xlabel={perte du portefeuille $x$ (\% du notionnel)},ylabel={$\P(L_5 > x)$},
  tick label style={font=\small},label style={font=\small},
  xmin=0,xmax=30,ymin=1e-4,ymax=1,grid=major,grid style={gray!25},
  legend columns=2,legend style={font=\footnotesize,at={(0.5,-0.25)},anchor=north,draw=none,fill=none,/tikz/every even column/.append style={column sep=8pt}},
  table/col sep=comma]
\addplot[omDef,very thick] table[x=x,y=gauss]{figdata/rates-credit-risk/15-portfolio-credit/tails.csv};
\addplot[omThm,very thick,dashed] table[x=x,y=t4]{figdata/rates-credit-risk/15-portfolio-credit/tails.csv};
\legend{copule gaussienne,copule de Student ($\nu = 4$)}
\end{axis}
\end{tikzpicture}

\omcaption{Probabilité que la perte du portefeuille à cinq ans dépasse $x$ (échelle logarithmique), mêmes
marginales et corrélation d'actifs de 25~\%, simulée avec 400\,000 tirages des facteurs. La
\omterm{def:rc:portfolio-credit:t}{copule de Student} accorde plusieurs fois plus de probabilité aux pertes importantes. Données~: le
tutoriel du chapitre.}\label{fig:rc:portfolio-credit:tails}
\end{omfigure}

\section{Mai 2005 et 2008}

\begin{dated}{2026-09}{Mai 2005}\label{dat:rc:portfolio-credit:may2005}
Selon le examen de la Banque des règlements internationaux de juin 2005, la vente de protection sur la
tranche equity de l'iTraxx Europe couverte par une protection sur l'indice a rapporté 300 à
400 points de base au début de 2005~; en mai, la corrélation implicite des prix de la tranche
equity a fortement baissé, les spreads des equity se sont élargis (plus vite après les dégradations
de General Motors et de Ford), les dénouements ont fait baisser les spreads de la tranche mezzanine,
et à certains moments la mezzanine a évolué à l'inverse de l'indice.
\end{dated}

Les deux épisodes sollicitent des parties différentes du modèle. En mai 2005, le risque était
idiosyncratique~: l'effondrement d'un ou deux noms a frappé la tranche equity, qui absorbe les
premières pertes, bien plus qu'une couverture de la taille de l'indice, et la corrélation
implicite des prix equity a baissé~; le delta, calculé pour un mouvement parallèle à corrélation
constante, ne couvrait ni l'un ni l'autre. En 2008, le risque était systémique~: le facteur a bougé,
les tranches senior que la \omterm{def:rc:portfolio-credit:gauss}{copule gaussienne} avait évaluées comme presque sans risque ont perdu de
la valeur, et les corrélations implicites des tranches senior ont grimpé vers un.

\section{Tutoriel~: tranches par récursion}\label{tut:rc:portfolio-credit:tutorial}

\begin{tutorial}
\textbf{Objectif.} Évaluer une structure de tranches sur un portefeuille de 125 noms, en déduire des
\omterm{def:rc:portfolio-credit:compound}{corrélations composées}, comparer les copules et calculer les deltas. \textbf{État final~:} les résultats
des \cref{ex:rc:portfolio-credit:quotes,ex:rc:portfolio-credit:deltas,ex:rc:portfolio-credit:tails} et
des \cref{fig:rc:portfolio-credit:lossdist,fig:rc:portfolio-credit:mezz,fig:rc:portfolio-credit:deltas,fig:rc:portfolio-credit:tails}.
\begin{tutsteps}
\item \textbf{Distribution des pertes}~: \texttt{loss\_distribution(p, rho)} pour trois
corrélations.
\item \textbf{Cotations}~: \texttt{quotes()} à partir du skew de corrélation de base.
\item \textbf{Composée}~: \texttt{compound\_table()}~; tracer \texttt{mezz\_curve()}.
\item \textbf{Queues et deltas}~: \texttt{tail\_table()}, \texttt{deltas()}~;
\texttt{fig\_rc\_portcredit.py} produit les graphiques.
\end{tutsteps}
\textbf{Pour aller plus loin.} Donner à dix noms un spread de 300 points de base et au reste 40,
au même niveau d'indice, puis comparer le prix de la tranche equity~; porter la cotation de la mezzanine
à 400 et observer la disparition de la \omterm{def:rc:portfolio-credit:compound}{corrélation composée}.
\end{tutorial}

\section{Construction~: portefeuille de crédit}\label{bld:rc:portfolio-credit:portcredit}

\begin{build}
\bfield{Rôle} Les modèles de portefeuille de la société~: distributions des pertes des portefeuilles de crédit
pour les tranches, pour les portefeuilles de contreparties de la partie~III, et pour le capital de crédit
de la partie~IV.
\bfield{Interface} \texttt{Pool(hazards, recovery)}~; \texttt{conditional\_pd},
\texttt{default\_count\_given\_z}, \texttt{loss\_distribution}~; \texttt{expected\_tranche\_losses},
\texttt{base\_el}, \texttt{legs\_from\_el}, \texttt{tranche\_price}~;
\texttt{compound\_correlations}~; \texttt{t\_cdf}, \texttt{t\_inv}, \texttt{simulate\_losses}~;
\texttt{default\_correlation}.
\bfield{Règles} Intensités de défaut plates par nom, notionnels égaux, un seul taux de recouvrement~; primes
trimestrielles sur le notionnel en cours attendu~; quadrature de Gauss--Hermite à 64~points~; le
\texttt{firm\_tranche} du livre~2 est la limite des grands portefeuilles.
\bfield{Tests d'acceptation} \texttt{code/firm/portcredit/tests/}~: la récursion sans
corrélation est binomiale~; un grand portefeuille converge vers la limite du livre~2~; la base et la composée
coïncident sur un skew plat~; l'indice ne dépend pas de la corrélation~; quantiles de Student~; pertes moyennes simulées.
\bfield{Pour aller plus loin} Notionnels et recouvrements inégaux (pertes sur une grille)~; recouvrement
stochastique~; les évaluateurs à chargement factoriel aléatoire et de Student par quadrature~; pools
sur mesure ramenés aux corrélations de base de l'indice.
\end{build}

\begin{omsources}
\item D.\ X.\ Li, ``On default correlation: a copula function approach'', \emph{Journal of
Fixed Income} 9(4), 2000, 43--54.
\item O.\ A.\ Vasicek, \emph{Finance, Economics and Mathematics}, Wiley, 2015, chapter
``Loan portfolio value''.
\item J.\ Hull and A.\ White, ``Valuation of a CDO and an $n$-th to default CDS without
Monte Carlo simulation'', \emph{Journal of Derivatives} 12(2), 2004, 8--23.
\item Bank for International Settlements, \emph{Quarterly Review}, June 2005.
\item Financial Crisis Inquiry Commission, \emph{The Financial Crisis Inquiry Report},
2011.
\item L.\ Andersen, J.\ Sidenius and S.\ Basu, ``All your hedges in one basket'', \emph{Risk},
November 2003.
\item Swiss Federal Banking Commission, \emph{Subprime crisis: SFBC investigation into the causes of the
write-downs of UBS AG}, 30 September 2008.
\end{omsources}

\section{Exercices}

\begin{exercise}[$\star$]\label{exo:rc:portfolio-credit:1}
Un portefeuille de 125 noms affiche chacun une \omterm{def:rc:reduced-form-credit:pd}{probabilité de défaut} à cinq ans de 4.08~\%. Combien de
défauts attend-on, et cette espérance dépend-elle de la corrélation~?
\end{exercise}

\begin{exercise}[$\star$]\label{exo:rc:portfolio-credit:2}
Quelles tranches sont longues de corrélation et lesquelles courtes, et pourquoi la mezzanine l'est-elle à la fois~?
\end{exercise}

\begin{exercise}[$\star$]\label{exo:rc:portfolio-credit:3}
Avec un recouvrement de 40~\%, combien de défauts parmi 125 noms effacent la tranche 0--3~\%, et combien
atteignent la \omterm{def:rc:portfolio-credit:ss}{tranche super senior} à 22~\%~?
\end{exercise}

\begin{exercise}[$\star\star$]\label{exo:rc:portfolio-credit:4}
Expliquer pourquoi une corrélation d'actifs de 25~\% ne donne qu'une \omterm{def:rc:portfolio-credit:dc}{corrélation de défaut} de 7~\%.
\end{exercise}

\begin{exercise}[$\star\star$]\label{exo:rc:portfolio-credit:5}
En utilisant les deltas de l'\cref{ex:rc:portfolio-credit:deltas}, quelle protection sur l'indice couvre
50~millions d'euros de protection vendue sur les tranches 6--9~\%~?
\end{exercise}

\begin{exercise}[$\star\star$]\label{exo:rc:portfolio-credit:6}
Pourquoi le marché cote-t-il des corrélations de base plutôt que des \omterm{def:rc:portfolio-credit:compound}{corrélations composées}, et
qu'a-t-elle de problématique, la corrélation de base~?
\end{exercise}

\begin{exercise}[$\star\star\star$]\label{exo:rc:portfolio-credit:7}
\emph{Programmation.} Calculer la probabilité que le portefeuille perde plus de 24~\% en cinq ans
sous les deux copules, pour une corrélation d'actifs de 25~\%, et le rapport de ces deux probabilités.
\end{exercise}

\begin{exercise}[$\star\star\star$]\label{exo:rc:portfolio-credit:8}
\emph{Trouver la faille.} «~Notre position super senior affiche une perte attendue de quelques points de base
et une notation AAA. Elle n'exige ni couverture ni presque de fonds propres.~»
\end{exercise}

\section{Problème~: long sur l'equity, short sur la mezzanine}

\begin{problem}[{Problème du week-end --- le trade de corrélation de mai 2005}]\label{pb:rc:portfolio-credit:1}
Un fonds vend 10~millions d'euros de protection 0--3~\% à cinq ans aux cotations de la
\cref{ex:rc:portfolio-credit:quotes} (prime initiale de 37.21~\% plus 500 de \omterm{def:rc:mortgage-modelling:cc}{coupon courant}) et achète une protection 3--6~\%
à 195 points de base, au montant neutre en delta. Puis le spread d'un émetteur passe
de 50 à 1\,000 points de base, et la corrélation de base de l'equity baisse de 15~\% à 12~\%,
les autres corrélations de base restant inchangées.

\medskip\noindent\textbf{Partie I --- La position.}
\begin{enumerate}
\item Donner les deltas des deux tranches et le notionnel de la mezzanine.
\item Que reçoit le fonds à l'initiation, et que gagne-t-il tant que rien
ne se passe~?
\item Quelle est son exposition à un élargissement parallèle, au défaut d'un émetteur
et à la corrélation~?
\item Pourquoi le trade a-t-il été appelé long corrélation~?
\item Quels risques le delta ne mesurait-il pas~?
\end{enumerate}

\medskip\noindent\textbf{Partie II --- Un émetteur s'effondre.}
\begin{enumerate}[resume]
\item Donner la prime initiale de l'equity et le spread de la mezzanine après le choc.
\item Donner le P\&L de chaque jambe et le total.
\item Pourquoi la couverture ne compense-t-elle pas la perte de l'equity~?
\item Que coûterait le défaut de cet émetteur (recouvrement de 40~\%) à la tranche
equity directement, en notionnel~?
\item Une couverture par l'indice aurait-elle mieux fonctionné~?
\end{enumerate}

\medskip\noindent\textbf{Partie III --- La corrélation baisse.}
\begin{enumerate}[resume]
\item Donner le P\&L de chaque jambe lorsque seule la corrélation de base de l'equity baisse à 12~\%.
\item Pourquoi le spread de la mezzanine baisse-t-il~?
\item Donner les P\&L des deux jambes et le total lorsque les deux événements se produisent.
\item Quelle jambe le fonds attendait-il pour le protéger, et que lui est-il arrivé~?
\item Comment le dénouement de positions similaires fait-il encore évoluer les prix~?
\end{enumerate}

\medskip\noindent\textbf{Partie IV --- Jugement.}
\begin{enumerate}[resume]
\item Quelle est la corrélation dont le fonds était long, dans le modèle et sur le marché~?
\item Comment limiter un tel portefeuille~?
\item Pourquoi les mêmes modèles n'ont-ils donné aucun signal avant les pertes sur les tranches senior de 2008~?
\item Énoncer le \emph{résultat nommé}~: les P\&L des deux jambes et le total
lorsqu'un émetteur s'effondre et que la corrélation de l'equity baisse.
\item En une phrase~: sur quoi une position de tranches couverte en delta parie-t-elle encore~?
\end{enumerate}
\end{problem}

\section{Questions d'entretien}

\begin{interviewq}[$\star$ \iqroles{chercheur, trader}]\label{iq:rc:portfolio-credit:1}
Écrivez le modèle de \omterm{def:rc:portfolio-credit:gauss}{copule gaussienne} à un facteur et la \omterm{def:rc:reduced-form-credit:pd}{probabilité de défaut} conditionnelle.
\end{interviewq}

\begin{interviewq}[$\star\star$ \iqroles{chercheur, développeur}]\label{iq:rc:portfolio-credit:2}
Comment calculez-vous la distribution des pertes d'un portefeuille fini d'émetteurs hétérogènes
sans recourir à Monte Carlo~?
\end{interviewq}

\begin{interviewq}[$\star\star$ \iqroles{trader}]\label{iq:rc:portfolio-credit:3}
\omterm{def:rc:portfolio-credit:compound}{Corrélation composée} contre corrélation de base~: qu'est-ce que chacune, et que va-t-il mal avec chacune~?
\end{interviewq}

\begin{interviewq}[$\star\star$ \iqroles{risque}]\label{iq:rc:portfolio-credit:4}
Comment gériez-vous le risque d'un portefeuille de \omterm{def:rc:portfolio-credit:ss}{tranches super senior}~?
\end{interviewq}

\begin{interviewq}[$\star\star\star$ \iqroles{chercheur}]\label{iq:rc:portfolio-credit:5}
Qu'est-ce que la dépendance de queue, quelles copules la présentent, et pourquoi
importe-t-elle pour les tranches senior~?
\end{interviewq}

\begin{interviewq}[$\star\star\star$ \iqroles{trader, risque}]\label{iq:rc:portfolio-credit:6}
Expliquez comment une position couverte en delta, longue equity et courte mezzanine, a perdu de l'argent en mai 2005.
\end{interviewq}
